% =====================================================================
% BMED365: Generative AI og large language models
% Beamer-presentasjon - Topic List G01-G14
% =====================================================================
\documentclass[aspectratio=169, 10pt]{beamer}

% =====================================================================
% PACKAGES
% =====================================================================
\usepackage[utf8]{inputenc}
\usepackage[T1]{fontenc}
\usepackage[notherwisek]{babel}
\usepackage{graphicx}
\usepackage{tikz}
\usetikzlibrary{shapes.geometric, arrows, positioning, calc, decorations.pathreplacing}
\usepackage{booktabs}
\usepackage{amsmath}
\usepackage{fontawesome5}
\usepackage{hyperref}
\hypersetup{colorlinks=true, linkcolor=blue, urlcolor=blue}

% =====================================================================
% THEME AND COLORS
% =====================================================================
\usetheme{Madrid}
\usecolortheme{beaver}

% =====================================================================
% TITLE INFO
% =====================================================================
\title{Generative AI and Large Language Models}
\subtitle{BMED365: Topic List G01--G20}
\author{BMED365}
\date{Spring 2026}

% =====================================================================
% DOCUMENT
% =====================================================================
\begin{document}

% Title page
\begin{frame}
    \titlepage
\end{frame}

% Table of Contents
\begin{frame}{Overview}
    \tableofcontents
\end{frame}

% =====================================================================
% SECTION: Introduction to Generative AI
% =====================================================================
\section{Introduction to Generative AI}

\begin{frame}{G01: Define generative AI and distinguish from discriminative AI}
    \begin{columns}[T]
        \begin{column}{0.48\textwidth}
            \textbf{Diskriminativ AI:}
            \begin{itemize}
                \item Lærer å \textbf{classify} or distinguish between kategorier
                \item Modorer $P(y|x)$ -- sannsynlighet for klasse $y$ gitt input $x$
                \item Examples: ``Er dette kreft?'', ``Hvilken disease?''
                \item CNN for bildeclassification
            \end{itemize}
        \end{column}
        \begin{column}{0.48\textwidth}
            \textbf{Generative AI:}
            \begin{itemize}
                \item Lærer å \textbf{skape} nytt innhold
                \item Modorer $P(x)$ or $P(x|y)$ -- dataadvantageingen selv
                \item Examples: Tekst, bilder, kode, musikk
                \item LLM, diffusjonsmodels
            \end{itemize}
        \end{column}
    \end{columns}

    \vspace{0.2cm}
    \begin{block}{\footnotesize Key difference}
        \footnotesize
        \textbf{Diskriminativ:} Input $x$ $\rightarrow$ Kategori/Label $y$ \quad | \quad
        \textbf{Generativ:} Prompt $\rightarrow$ Nytt innhold
    \end{block}

    \begin{alertblock}{\footnotesize Medical application}
        \footnotesize Generative AI: Oppsummere journaler, generere rapporter, svare på spørsgoal, syntetisere training data
    \end{alertblock}
\end{frame}

% =====================================================================
% SECTION: Transformer Architecture
% =====================================================================
\section{Transformer Architecture}

\begin{frame}{G02: Explain \href{https://en.wikipedia.org/wiki/Attention_(machine_learning)}{self-attention}-mekanismen at a conceptual level}
    \textbf{Self-attention: ``What is relevant for hva?''}

    \vspace{0.2cm}
    \textbf{Intuisjon:}
    \begin{itemize}
        \item For hvert ord: Se på \textbf{alle} andre ord i setningen
        \item Beregn en \textbf{vektet sum} basert på relevans
        \item Fanger opp avhengigheter over lange avstander
    \end{itemize}

    \vspace{0.2cm}
    \textbf{Example:}
    \begin{center}
        \textit{``Pasienten tok \textbf{medisinen}, og \textbf{den} hjalp mot smertene.''}
    \end{center}
    \begin{itemize}
        \item Self-attention kobler ``den'' til ``medisinen'' (ikke ``smertene'')
    \end{itemize}

    \vspace{0.1cm}
    \begin{block}{\footnotesize \href{https://jalammar.github.io/illustrated-transformer/\#self-attention-at-a-high-level}{Query, Key, Value (QKV)}}
        \footnotesize Hvert token genererer tre vekcountser: \\
        \href{https://en.wikipedia.org/wiki/Attention_(machine_learning)\#Queries,_keys,_and_values}{\textbf{Query}} ``Hva leter jeg etter?'' |
        \href{https://en.wikipedia.org/wiki/Attention_(machine_learning)\#Queries,_keys,_and_values}{\textbf{Key}} ``Hva kan jeg tilby?'' |
        \href{https://en.wikipedia.org/wiki/Attention_(machine_learning)\#Queries,_keys,_and_values}{\textbf{Value}} ``What is innholdet mitt?''
    \end{block}
\end{frame}

\begin{frame}{G03: Describe the transformer architecture (\href{https://en.wikipedia.org/wiki/Transformer_(deep_learning_architecture)\#Encoder}{encoder}-\href{https://en.wikipedia.org/wiki/Transformer_(deep_learning_architecture)\#Decoder}{decoder})}
    \textbf{Transformer Architecture (\href{https://arxiv.org/abs/1706.03762}{Vaswani et al., 2017}):}

    \begin{columns}[T]
        \begin{column}{0.55\textwidth}
            \textbf{Encoder:}
            \begin{itemize}
                \item Processes input-sekvens
                \item Bygger kontekstrik representasjon
                \item Brukes i: \href{https://arxiv.org/abs/1810.04805}{BERT}, \href{https://en.wikipedia.org/wiki/Word_embedding}{embeddings}
            \end{itemize}

            \vspace{0.2cm}
            \textbf{Decoder:}
            \begin{itemize}
                \item Genererer output sekvensielt
                \item Bruker \href{https://en.wikipedia.org/wiki/Attention_(machine_learning)\#Causal_attention}{maskert self-attention}
                \item Brukes i: GPT, Claude, Gemini
            \end{itemize}

            \vspace{0.2cm}
            \textbf{Full encoder-decoder:}
            \begin{itemize}
                \item Oversettelse, oppsummering
                \item \href{https://arxiv.org/abs/1910.10683}{T5}, \href{https://arxiv.org/abs/2210.11416}{Flan-T5}
            \end{itemize}
        \end{column}
        \begin{column}{0.42\textwidth}
            \textbf{Nøkkelkomponenter:}
            \begin{enumerate}
                \item \href{https://en.wikipedia.org/wiki/Attention_(machine_learning)\#Multi-head_attention}{Multi-head self-attention}
                \item \href{https://en.wikipedia.org/wiki/Feedforward_neural_network}{Feed-forward nettverk}
                \item \href{https://en.wikipedia.org/wiki/Layer_normalization}{Layer normalization}
                \item \href{https://en.wikipedia.org/wiki/Residual_neural_network}{Residual connections}
                \item \href{https://en.wikipedia.org/wiki/Positional_encoding}{Positional encoding}
            \end{enumerate}

            \vspace{0.2cm}
            \begin{block}{\footnotesize GPT = ``Decoder-only''}
                \footnotesize Scountse language models (LLM) bruker typisk kun decoder-delen for tekstgenerering.
            \end{block}
        \end{column}
    \end{columns}
\end{frame}

% =====================================================================
% SECTION: LLM Fundamentals
% =====================================================================
\section{LLM Fundamentals}

\begin{frame}{G04: Explain hva \href{https://en.wikipedia.org/wiki/Lexical_analysis\#Token}{tokens} er og how \href{https://en.wikipedia.org/wiki/Tokenization_(lexical_analysis)}{tokenisering} fungerer}
    \textbf{Tokens = tekstens ``atomer''}
    \begin{itemize}
        \item LLM-er leser ikke bokstaver or ord -- de leser \textbf{tokens}
        \item Token $\approx$ delord, ca. 4 tegn per token (for engelsk/notherwisek)
    \end{itemize}

    \vspace{0.15cm}
    \textbf{Tokenisering -- eksempel:}
    \begin{center}
        \footnotesize
        \texttt{``Pasienten har diabetes''} $\rightarrow$ \texttt{[``Pas'', ``ient'', ``en'', `` har'', `` diab'', ``etes'']}
    \end{center}

    \vspace{0.1cm}
    \begin{columns}[T]
        \begin{column}{0.48\textwidth}
            \textbf{\href{https://en.wikipedia.org/wiki/Byte_pair_encoding}{BPE} og språkfotherwisekjor:}
            \footnotesize
            \begin{itemize}
                \item Bygger vokabular iterativt
                \item GPT-4: ca. 100 000 tokens
                \item \textbf{Kinesisk:} Hvert tegn $\approx$ 1--2 tokens
                \item Forenklet/tradisjonell: Ulike tokens
            \end{itemize}
        \end{column}
        \begin{column}{0.48\textwidth}
            \textbf{Important å vite:}
            \footnotesize
            \begin{itemize}
                \item Lange/sjeldne ord = flere tokens
                \item Påvirker kostnad og hastighet
                \item \href{https://github.com/openai/tiktoken}{\texttt{tiktoken}} (OpenAI) for telling
                \item Ulike models = ulike tokenizers
            \end{itemize}
        \end{column}
    \end{columns}

    \begin{block}{\footnotesize Why tokens?}
        \footnotesize Mer effektivt enn bokstaver (for kort) or hele ord (vokabular for scountst).
    \end{block}
\end{frame}

\begin{frame}{G05: Understand the concept of \href{https://en.wikipedia.org/wiki/Context_window}{kontekstvindu} (context window)}
    \textbf{Kontekstvindu = modellens ``arbeidsminne''}
    \begin{itemize}
        \item Maksimalt antall tokens modellen kan ``se'' samtidig
        \item Inkluderer både \textbf{input} (prompt) og \textbf{output} (svar)
    \end{itemize}

    \vspace{0.2cm}
    \begin{columns}[T]
        \begin{column}{0.55\textwidth}
            \textbf{Examples på kontekstvinduer:}
            \footnotesize
            \begin{tabular}{lrl}
                \toprule
                Modell & Kontekst & $\approx$ Sider \\
                \midrule
                GPT-4o & 128K & $\sim$150 sider \\
                GPT-5 & 256K+ & $\sim$300 sider \\
                Claude Opus 4.5 & 200K & $\sim$250 sider \\
                Gemini 3 Pro & 2M+ & Hele bøker! \\
                \bottomrule
            \end{tabular}
        \end{column}
        \begin{column}{0.42\textwidth}
            \textbf{Input vs. output:}
            \footnotesize
            \begin{itemize}
                \item \textbf{Input-kontekst:} Hvor mye modellen kan lese
                \item \textbf{Output-kontekst:} Hvor langt svar den kan gi
                \item Output ofte mindre (4K--32K tokens)
            \end{itemize}
        \end{column}
    \end{columns}

    \vspace{0.1cm}
    \begin{alertblock}{\footnotesize Limitations}
        \footnotesize
        For lange tekster må kuttes. Modellen ``glemmer'' utenfor vinduet. Større vindu = dyrere/tregere.
    \end{alertblock}
\end{frame}

\begin{frame}{G06: Explain hva \href{https://en.wikipedia.org/wiki/Softmax_function\#Softmax_temperature}{temperatur} betyr i tekstgenerering}
    \textbf{Temperatur = kontroll over ``kreativitet''}

    \vspace{0.15cm}
    \textbf{Teknisk:}
    \begin{itemize}
        \item Skalerer logits før softmax: $p_i = \frac{e^{z_i/T}}{\sum_j e^{z_j/T}}$
        \item \textbf{Effekt:} Lav $T$ $\rightarrow$ skarpere advantageing (høyeste logit dominerer); Høy $T$ $\rightarrow$ flatere advantageing (mer tilfeldighet)
    \end{itemize}

    \vspace{0.2cm}
    \begin{columns}[T]
        \begin{column}{0.48\textwidth}
            \textbf{Lav temperatur (0.0--0.3):}
            \begin{itemize}
                \item Mer deterministisk
                \item Velger mest sannsynlige tokens
                \item Konsistent, ``trygt''
                \item \faCheck~Medical info, fakta
            \end{itemize}
        \end{column}
        \begin{column}{0.48\textwidth}
            \textbf{Høy temperatur (0.7--1.0+):}
            \begin{itemize}
                \item Mer tilfeldig/kreativ
                \item Flere overraskende valg
                \item Mer variasjon
                \item \faCheck~Kreativ skriving, brainscountsming
            \end{itemize}
        \end{column}
    \end{columns}

    \vspace{0.15cm}
    \begin{block}{\footnotesize Recommendation for medicine}
        \footnotesize Bruk \textbf{lav temperatur} (0.0--0.3) for faktabaserte taskr. Høyere temperatur øker risiko for hallusinering.
    \end{block}
\end{frame}

% =====================================================================
% SECTION: Prompt Engineering
% =====================================================================
\section{Prompt Engineering}

\begin{frame}{G07: Apply \href{https://en.wikipedia.org/wiki/Zero-shot_learning}{zero-shot} prompting}
    \textbf{Zero-shot = ingen eksempler} -- be modellen løse taskn \textbf{without vise eksempler}

    \vspace{0.1cm}
    \begin{columns}[T]
        \begin{column}{0.55\textwidth}
            \textbf{Examples på zero-shot prompts:}
            \begin{block}{\footnotesize Zero-shot prompts}
                \scriptsize
                \begin{enumerate}
                    \item \texttt{Klassifiser symptom som akutt/kronisk: ``Hodepine i 3 mnd''}
                    \item \texttt{Oppsummer journalnotatet i 3 setninger.}
                    \item \texttt{Oversett til patientvennlig språk: ``Bilateral pneumoni''}
                    \item \texttt{Ekstraher alle medikamenter fra denne teksten.}
                    \item \texttt{Er denne lab-valueen unormal? Hb 8.2 g/dL}
                \end{enumerate}
            \end{block}
        \end{column}
        \begin{column}{0.42\textwidth}
            \textbf{Når fungerer zero-shot?}
            \footnotesize
            \begin{itemize}
                \item Modellen har sett lignende taskr
                \item Oppgaven er tydelig definert
            \end{itemize}

            \vspace{0.1cm}
            \begin{alertblock}{\footnotesize Limitations}
                \scriptsize
                For spesialiserte taskr kan zero-shot gi dårlige resulter. Da trengs \textbf{few-shot}.
            \end{alertblock}
        \end{column}
    \end{columns}
\end{frame}

\begin{frame}{G08: Apply \href{https://en.wikipedia.org/wiki/Few-shot_learning}{few-shot prompting}}
    \textbf{Few-shot = noen eksempler}
    
    \begin{itemize}
        \item Gi modellen \textbf{2--5 eksempler} på ønsket input-output
        \item Modellen lærer mønsteret ``in-context''
    \end{itemize}
    
    \vspace{0.3cm}
    \begin{block}{Few-shot prompt}
        \small
        \texttt{Klassifiser symptomvarighet:} \\
        \texttt{Symptom: ``Brystsmerter i 30 minutter'' → Akutt} \\
        \texttt{Symptom: ``Ryggsmerter i 2 år'' → Kronisk} \\
        \texttt{Symptom: ``Feber siden i går'' → Akutt} \\
        \texttt{---} \\
        \texttt{Symptom: ``Hodepine i 3 måneder'' →}
    \end{block}
    
    \vspace{0.3cm}
    \textbf{Tip for effektiv few-shot:}
    \begin{itemize}
        \item Velg \textbf{representative} eksempler
        \item Vis \textbf{variasjonen} i mulige outputs
        \item Hold eksemplene \textbf{konsistente} i format
        \item Start med 3--5 eksempler, juster ved behov
    \end{itemize}
\end{frame}

\begin{frame}{G09: Apply \href{https://en.wikipedia.org/wiki/Prompt_engineering\#Chain-of-thought}{Chain-of-Thought (CoT)} prompting}
    \textbf{Chain-of-Thought = vis resonnementet}
    \begin{itemize}
        \item Be modellen \textbf{tenke steg-for-steg}
        \item Forbedrer ytelse på complexe resonneringstaskr
    \end{itemize}

    \vspace{0.2cm}
    \begin{columns}[T]
        \begin{column}{0.48\textwidth}
            \begin{alertblock}{\footnotesize Uten CoT}
                \footnotesize
                \texttt{Bør patienten henvises?} \\
                \texttt{→ ``Ja''}
            \end{alertblock}
        \end{column}
        \begin{column}{0.48\textwidth}
            \begin{block}{\footnotesize Med CoT}
                \footnotesize
                \texttt{Tenk steg-for-steg:} \\
                \texttt{1. Symptomer: X, Y, Z} \\
                \texttt{2. Alvorlighetsgrad: Moderat} \\
                \texttt{→ Konklusjon: Bør henvises}
            \end{block}
        \end{column}
    \end{columns}

    \vspace{0.15cm}
    \textbf{Variants:}
    \begin{itemize}
        \item \textbf{Zero-shot CoT:} ``Let's think step by step''
        \item \textbf{Few-shot CoT:} Vis eksempler med resonnement
    \end{itemize}

    \begin{block}{\footnotesize Built-in CoT i nyere LLM-er}
        \footnotesize ChatGPT o1/o3 (Thinking), Claude Opus 4.5 (Extended Thinking), Gemini 3 (Deep Think) har CoT innebygd.
    \end{block}
\end{frame}

\begin{frame}{G10: Describe god praksis for \href{https://en.wikipedia.org/wiki/Prompt_engineering\#System_prompts}{systemprompts}}
    \textbf{Systemprompt = instruksjon som setter kontekst}

    \vspace{0.2cm}
    \textbf{Elementer i et godt systemprompt:}
    \begin{enumerate}
        \item \textbf{Rolle:} ``Du er en medical assistent...''
        \item \textbf{Kontekst:} ``Brukeren er helsepersonell...''
        \item \textbf{Oppførsel:} ``Svar konsist og presist...''
        \item \textbf{Limitations:} ``Ikke gi diagnosisr...''
        \item \textbf{Format:} ``Svar på notherwisek, bruk punktlister...''
    \end{enumerate}

    \vspace{0.15cm}
    \begin{block}{\footnotesize Example på systemprompt}
        \scriptsize
        \texttt{Du er en medical AI-assistent for helsepersonell. Svar på notherwisek. Vær presis og evidensbasert. Henvis til sourcer når mulig. Gi aldri definitive diagnosisr -- foreslå differensialdiagnosisr. Ved usikkerhet, si fra.}
    \end{block}

    \begin{alertblock}{\footnotesize Tip}
        \footnotesize Test og iterer på systemprompts! Små endringer kan gi scounts effekt.
    \end{alertblock}
\end{frame}

% =====================================================================
% SECTION: Context Engineering
% =====================================================================
\section{Context Engineering}

\begin{frame}{G15: Fra \href{https://en.wikipedia.org/wiki/Prompt_engineering}{Prompt Engineering} til \href{https://www.anthropic.com/research/building-effective-agents}{Context Engineering}}
    \textbf{\href{https://en.wikipedia.org/wiki/Prompt_engineering}{Prompt Engineering} vs. \href{https://www.anthropic.com/research/building-effective-agents}{Context Engineering}:}
    \begin{itemize}
        \item \textbf{\href{https://en.wikipedia.org/wiki/Prompt_engineering}{Prompt engineering}:} Fokus på å utforme gode spørsgoal
        \item \textbf{\href{https://www.anthropic.com/research/building-effective-agents}{Context engineering}:} Designe \textit{hele kontekstvinduet} strategisk
    \end{itemize}
    
    \vspace{0.2cm}
    \textbf{Komponenter i context engineering:}
    \footnotesize
    \begin{tabular}{lll}
        \toprule
        \textbf{Komponent} & \textbf{Description} & \textbf{Medical example} \\
        \midrule
        System prompt & Rolle og instruksjoner & ``Clinical beslutningsstøtte'' \\
        Bruker-prompt & Spesifikt spørsgoal & ``Vurder patientens symptoms'' \\
        Verktøy & Tilgjengelige handlinger & Lab-søk, ICD-10, Felleskatalogen \\
        RAG-kontekst & Hentet informasjon & Retningslinjer, journal, artikler \\
        Samtalehiscountsikk & Tidligere dialog & Tidligere spørsgoal og svar \\
        Strukturerte data & Formattert info & JSON med vitale tegn \\
        \bottomrule
    \end{tabular}
    
    \normalsize
    \vspace{0.15cm}
    \begin{block}{\footnotesize Why important i medisin?}
        \footnotesize Rik kontekst $\rightarrow$ sikrere, mer relevante og personaliserte svar
    \end{block}
\end{frame}

\begin{frame}{G16--G17: Model-specific prompting techniques}
    \textbf{Tre ledende LLM-leverandører med ulike strengthr:}
    
    \vspace{0.2cm}
    \begin{columns}[T]
        \begin{column}{0.32\textwidth}
            \textbf{\href{https://openai.com/gpt-5}{GPT-5.2} (OpenAI):}
            \footnotesize
            \begin{itemize}
                \item Sterkest på strukturert output
                \item Beste verktøyintegrasjon
                \item God instruksjonsetterlevelse
            \end{itemize}
            \textbf{Godt egnet for:}
            \begin{itemize}
                \item JSON-ekstraksjon
                \item API-integrasjoner
                \item Agentiske systemer
            \end{itemize}
        \end{column}
        \begin{column}{0.32\textwidth}
            \textbf{\href{https://anthropic.com/claude}{Claude Opus 4.5}:}
            \footnotesize
            \begin{itemize}
                \item 200K tokens kontekst
                \item Nyansert resonnering
                \item Constitutional AI
            \end{itemize}
            \textbf{Godt egnet for:}
            \begin{itemize}
                \item Etiske dilemmaer
                \item Lang journalanalyse
                \item Pasientkommunikasjon
            \end{itemize}
        \end{column}
        \begin{column}{0.32\textwidth}
            \textbf{\href{https://deepmind.google/gemini}{Gemini 3} (Google):}
            \footnotesize
            \begin{itemize}
                \item Multimodal (bilde+tekst)
                \item Sanntids websøk
                \item Vitenskapelig strength
            \end{itemize}
            \textbf{Godt egnet for:}
            \begin{itemize}
                \item Image Analysis
                \item Litteratursøk
                \item Beregninger
            \end{itemize}
        \end{column}
    \end{columns}
    
    \vspace{0.15cm}
    \begin{block}{\footnotesize Lab 3 Notebook 04}
        \footnotesize Inneholder detaljerte prompting-guider for hver modell + sammenligningseksperimenter
    \end{block}
\end{frame}

\begin{frame}{G18: Handle uncertainty and hallucination risk}
    \textbf{Prompt-mønster for medical usikkerhetshåndtering:}
    
    \vspace{0.2cm}
    \begin{block}{\footnotesize Anti-hallucination prompt (GPT-5.2 stil)}
        \scriptsize
        \texttt{<usikkerhet\_og\_tvetydighet>} \\
        \texttt{- Ved tvetydige symptoms: Presenter differensialdiagnosisr med sannsynlighet} \\
        \texttt{- Oppgi tydelig kunnskapens limitationer og dato for update} \\
        \texttt{- Henvis alltid til gjeldende retningslinjer og clinical vurdering} \\
        \texttt{- Ved usikkerhet: Eshallér til høyere hastegrad} \\
        \texttt{</usikkerhet\_og\_tvetydighet>}
    \end{block}
    
    \vspace{0.15cm}
    \textbf{Selvsjekk for høyrisiko-taskr:}
    \footnotesize
    \begin{enumerate}
        \item Skann eget svar for usagte antakelser
        \item Verifiser at tall og påstander er forankret i kontekst
        \item Unngå for sterk språkbruk (``alltid'', ``garantert'')
        \item Kvalifiser konklusjoner med usikkerhetsgrad
    \end{enumerate}
    
    \normalsize
    \begin{alertblock}{\footnotesize Rule of thumb}
        \footnotesize AI = beslutningsstøtte, \textbf{ikke} beslutningstaker. Alltid clinical validation!
    \end{alertblock}
\end{frame}

\begin{frame}{G19--G20: Structured extraction and model comparison}
    \begin{columns}[T]
        \begin{column}{0.48\textwidth}
            \textbf{G19: Strukturert ekstraksjon:}
            \footnotesize
            \begin{block}{\scriptsize JSON-skjema for epikrise}
                \scriptsize
                \texttt{\{} \\
                \texttt{~~"diagnosis": \{} \\
                \texttt{~~~~"tekst": string,} \\
                \texttt{~~~~"icd10": string | null} \\
                \texttt{~~\},} \\
                \texttt{~~"drugs": [\{} \\
                \texttt{~~~~"navn": string,} \\
                \texttt{~~~~"dose": string} \\
                \texttt{~~\}]} \\
                \texttt{\}}
            \end{block}
            
            \footnotesize
            \textbf{Importante regler:}
            \begin{itemize}
                \item Bruk \texttt{null} for manglende felt
                \item Krev sourcereferanser
                \item Verifiser mot originaldokument
            \end{itemize}
        \end{column}
        \begin{column}{0.48\textwidth}
            \textbf{G20: Beslutningsmatrise:}
            \scriptsize
            \begin{tabular}{ll}
                \toprule
                \textbf{Oppgave} & \textbf{Modell} \\
                \midrule
                Strukturert output & GPT-5.2 \\
                Etiske vurderinger & Claude \\
                Image Analysis & Gemini \\
                Lang kontekst & Claude \\
                Sanntidssøk & Gemini \\
                Verktøy/API & GPT-5.2 \\
                \bottomrule
            \end{tabular}
            
            \vspace{0.2cm}
            \normalsize
            \textbf{Praktisk tips:}
            \footnotesize
            \begin{enumerate}
                \item Test samme prompt på 2--3 models
                \item Sammenlign mot evaluationskriterier
                \item Velg modell basert på task
                \item Dokumenter valg for etterprøvbarhet
            \end{enumerate}
        \end{column}
    \end{columns}
\end{frame}

% =====================================================================
% SECTION: Challenges with LLM
% =====================================================================
\section{Challenges with LLM}

\begin{frame}{G11: Define \href{https://en.wikipedia.org/wiki/Hallucination_(artificial_intelligence)}{hallusinering} and its implikasjoner}
    \textbf{Hallusinering = modellen ``finner på'' feil informasjon}
    \begin{itemize}
        \item LLM-er genererer tekst som \textit{høres} riktig ut, men er \textit{faktisk feil}
        \item Også kalt \href{https://en.wikipedia.org/wiki/Confabulation_(neural_networks)}{\textit{konfabulering}} -- begge termer brukes i litteraturen
    \end{itemize}

    \vspace{0.15cm}
    \begin{columns}[T]
        \begin{column}{0.55\textwidth}
            \textbf{Types of Hallucinations:}
            \footnotesize
            \begin{itemize}
                \item \textbf{Faktiske feil:} Feil statistikk, datoer, navn
                \item \textbf{Oppdiktede sourcer:} Referanser som ikke eksisterer
                \item \textbf{Logiske feil:} Inkonsistente resonnementer
            \end{itemize}
        \end{column}
        \begin{column}{0.42\textwidth}
            \textbf{Utvikling:}
            \footnotesize
            \begin{itemize}
                \item Nyere models hallusinerer \textit{mindre}
                \item Kan \textit{ikke} elimineres helt
                \item RAG og CoT reduserer risiko
            \end{itemize}
        \end{column}
    \end{columns}

    \vspace{0.1cm}
    \begin{alertblock}{\footnotesize Implikasjoner i medisin}
        \footnotesize
        Feil dosering can be livstruende. \textbf{Alltid verifiser} mot pålitelige sourcer! LLM = beslutningsstøtte, \textbf{ikke} erstatning for fagkunnskap.
    \end{alertblock}
\end{frame}

% =====================================================================
% SECTION: Models and Advanced Concepts
% =====================================================================
\section{Models and Advanced Concepts}

\begin{frame}{G12: Know about GPT-5, Claude, Gemini and their applications}
    \begin{center}
        \scriptsize
        \begin{tabular}{lp{3cm}p{3.5cm}p{2.5cm}}
            \toprule
            \textbf{Modell} & \textbf{Utvikler} & \textbf{Styrker} & \textbf{Application} \\
            \midrule
            \href{https://openai.com/gpt-5}{GPT-5/o3} & OpenAI & Multimodal, reasoning & ChatGPT, Copilot \\
            \href{https://www.anthropic.com/claude}{Claude 4} & Anthropic & 200K kontekst, ``trygg'' & Dokumentanalyse \\
            \href{https://deepmind.google/gemini}{Gemini 3} & Google & Integrert i Colab & Kodehjelp, søk \\
            \href{https://llama.meta.com/}{Llama 4} & Meta & Open source, lokalt & Fotherwisekning \\
            \bottomrule
        \end{tabular}
    \end{center}

    \vspace{0.1cm}
    \begin{columns}[T]
        \begin{column}{0.48\textwidth}
            \textbf{Lokale models:}
            \scriptsize
            \begin{itemize}
                \item \href{https://ollama.ai/}{Ollama}, \href{https://lmstudio.ai/}{LM Studio}
                \item \href{https://en.wikipedia.org/wiki/Knowledge_distillation}{Destillering}: Liten modell lærer fra scounts
                \item \href{https://en.wikipedia.org/wiki/Quantization_(signal_processing)\#Neural_network_quantization}{Kvantisering}: Redusert precision (16$\rightarrow$4 bit)
                \item GDPR-vennlig, offline
            \end{itemize}
        \end{column}
        \begin{column}{0.48\textwidth}
            \begin{block}{\scriptsize \href{https://chat.uib.no/}{chat.uib.no}}
                \scriptsize
                \textbf{Styrker:} GDPR-ok, ingen datalayersring, gratis for UiB\\
                \textbf{Svakheter:} Eldre modell, begrenset kontekst, ikke multimodal
            \end{block}
        \end{column}
    \end{columns}
\end{frame}

\begin{frame}{G13: Explain the concept of \href{https://en.wikipedia.org/wiki/Foundation_model}{foundation model (foundation model)}}
    \textbf{Foundation model = pretrent på enorme datamengder}
    \begin{itemize}
        \item Trent på terabytes av tekst/bilder uten spesifikk task
        \item Kan tilpasses (\textit{finetuned}) til mange fotherwisekjellige taskr
        \item Koster millioner å trene fra scratch
    \end{itemize}

    \vspace{0.2cm}
    \textbf{Karakteristikker:}
    \begin{enumerate}
        \item \textbf{Scounts scale:} Milliarder av parametre
        \item \textbf{Selvveiledet læring:} Ingen labels nødvendig
        \item \textbf{Emergente egenskaper:} Evner som ``dukker opp'' ved scale
        \item \textbf{Tilpasningsbar:} Finetuning, prompting, RAG
    \end{enumerate}

    \vspace{0.15cm}
    \begin{block}{\footnotesize Examples på foundation models (2025)}
        \scriptsize
        \textbf{Tekst:} \href{https://openai.com/gpt-5}{GPT-5}, \href{https://www.anthropic.com/claude}{Claude 4}, \href{https://llama.meta.com/}{Llama 4} |
        \textbf{Imager:} \href{https://openai.com/dall-e-3}{DALL-E 3}, \href{https://stability.ai/}{Stable Diffusion 3} |
        \textbf{Multimodal:} \href{https://deepmind.google/gemini}{Gemini 3}, \href{https://openai.com/sora}{Sora}
    \end{block}
\end{frame}

\begin{frame}{G14: Describe \href{https://en.wikipedia.org/wiki/Retrieval-augmented_generation}{RAG} (Retrieval-Augmented Generation)}
    \textbf{RAG = koble LLM til ekstern kunnskapsbase}

    \vspace{0.2cm}
    \begin{columns}[T]
        \begin{column}{0.48\textwidth}
            \textbf{Problemet RAG løser:}
            \begin{itemize}
                \item LLM har ``frossen'' kunnskap
                \item Kan hallusinere fakta
                \item Mangler domene-spesifikk info
            \end{itemize}
        \end{column}
        \begin{column}{0.48\textwidth}
            \textbf{How RAG works:}
            \begin{enumerate}
                \item \textbf{Spørring:} Bruker stiller spørsgoal
                \item \textbf{Retrieval:} Hent dokumenter
                \item \textbf{Augmentation:} Legg til prompt
                \item \textbf{Generation:} LLM svarer
            \end{enumerate}
        \end{column}
    \end{columns}

    \vspace{0.15cm}
    \begin{block}{\footnotesize Medical application}
        \scriptsize
        \begin{itemize}
            \item Koble LLM til retningslinjer, \href{https://www.felleskatalogen.no/}{Felleskatalogen}, \href{https://www.uptodate.com/}{UpToDate}
            \item Svare basert på patientjournal (med tilganger)
            \item Reduserer hallusinering by forankre svar i sourcer
            \item Clinical beslutningsstøtte med oppdatert evidens
            \item Søk i institusjons-spesifikke prosedyrer og protokoller
        \end{itemize}
    \end{block}
\end{frame}

% =====================================================================
% SUMMARY
% =====================================================================
\section*{Summary}

\begin{frame}{Summary: Generative AI og LLM}
    \textbf{Fundamental:}
    \begin{itemize}
        \item G01: Generativ vs. diskriminativ AI
        \item G02--G03: Self-attention og transformer-arkitektur
        \item G04--G06: Tokens, kontekstvindu, temperatur
    \end{itemize}

    \textbf{Prompt og context engineering:}
    \begin{itemize}
        \item G07--G10: Zero-shot, few-shot, CoT, systemprompts
        \item G15--G17: Context engineering, modellspesifikke teknikker
        \item G18--G20: Usikkerhetshåndtering, strukturert ekstraksjon
    \end{itemize}

    \textbf{Challengeer og avansert:}
    \begin{itemize}
        \item G11: Hallusinering -- critical in a medical context
        \item G12--G14: Modor (GPT-5.2, Claude, Gemini), RAG
    \end{itemize}

    \vspace{0.1cm}
    \begin{block}{\footnotesize Lab 3 Notebook 04}
        \footnotesize Detaljerte prompting-guider for GPT-5.2, Claude Opus 4.5 og Gemini 3 med sammenligningseksperimenter.
    \end{block}
\end{frame}

\end{document}
